\documentclass[10pt]{article}

\usepackage[utf8]{inputenc}

\usepackage[T1]{fontenc}

\usepackage{amsmath}

\usepackage{amsfonts}

\usepackage{amssymb}

\usepackage[version=4]{mhchem}

\usepackage{stmaryrd}

\usepackage{bbold}

\usepackage{mathrsfs}



\begin{document}

статистик играют важную роль в математич. статистике, в частности в задаче построения подобных критериев, обладающих Неймана структурой.



Лит.: [1] Л и н н и к Ю. В., Статистические задачи с мешающими параметрами, М., 1966; [2] J е м а н Э., Проверка статистических гипотез, пер. с англ., 2 изд., М., 1979.



РАСПРЕДЕЛЕНИИ СХОДИМОСТЬ - В $\begin{gathered}\text { ОСновном }\end{gathered}$ слабая сходимость и сходимость по вариации, определяөмые следующим образом. Последовательность распределений (вероятностных мер) $\left\{P_{n}\right\}$ на борелевских множествах метрич. пространства $S$ наз. с л а 6 о с х о я ш е й с я распределению $P$, если



\[

\lim _{n} \int_{s} f d P_{n}=\int_{s} f d P

\]



для любой действительной ограниченной непрерывной функции $f$ на $S$. Слабая сходимость является основным типом сходимости, рассматриваемым в теории вероятностей. Обозначают ее обычно знаком $\Rightarrow$. Следугщие условия равносильны слабой сходимости:



\begin{enumerate}

  \item соотношение (*) выполняется для любой ограниченной равномерно нешрерывной дейіствительной функции $f$;



  \item соотнопение (*) выполняется для любой ограниченной непрерывной $P$-почти всюду действительной функции f;



  \item $\varlimsup_{n} P_{n}(F) \leqslant P(F)$



\end{enumerate}



для любого замкнутого множества $F \subset S$;



\begin{enumerate}

  \setcounter{enumi}{3}

  \item $\frac{\lim }{n} P_{n}(G) \geqslant P(G)$

\end{enumerate}



для любого открытого множества $G \subset S$;



\begin{enumerate}

  \setcounter{enumi}{4}

  \item $\lim _{n} P_{n}(A)=P(A)$

\end{enumerate}



для любого борелевского множества $A \subset S$ такого, что $P(\partial A)=0$, где $\partial A$ - граница $A$;



\begin{enumerate}

  \setcounter{enumi}{5}

  \item $\lim _{n} p\left(P_{n}, P\right)=0$,

\end{enumerate}



где $p$ есть Леви- Прохорова метрика.



Пусть $U$ - замкнутый относительно псресетений класс подмножеств $S$ такой, что всякое открытое множество из $S$ есть конечное или счетное объединение мнояеств из $U$. Тогда если $\lim P_{n}(A)=P(A)$ при всех $A \in U$, то $P_{n} \Rightarrow P$. Если $S=\mathbb{R}^{k}$ и $F_{n}, F-$ функции распределения, отвечающие $P_{n}$ II $P$ соответственно, то $P_{n} \Rightarrow P$ тогда и только тогда, когда $F_{n}(x) \rightarrow F(x)$ в кажкдой точке $x$ непрерывности функции $F$.



Пусть пространство $S$ сегарабельно и of - класс ограниченных борелевских действительных функций на $S$. Для того чтобы $\int_{S} f d P_{n} \rightarrow \int_{S} f d P$ равномерно по $f \in \mathscr{F}$ для всякой последовательності $\left\{P_{n}\right\}$ такої, тто $P_{n} \Rightarrow P$, необходимо и достаточно, чтобы:



a) $\sup _{f \in F} \omega_{f}(S)<\infty$



б) $\left.\left.\lim _{\varepsilon \downarrow 0 f \in \mathcal{F}} \sup _{f\left(\left\{x: \omega_{f}\right.\right.}\left(S_{x}, \varepsilon\right)>\delta\right\}\right)=0, \forall \delta>0$,



где



\[

\omega_{f}(A)=\sup \{|f(x)-f(y)|, \quad x, y \in A\},

\]



и где $S_{x, \varepsilon}$ - открытый шар раднуса $\varepsilon$ с центром в $x$. Если класс Ғ образован индикаторами множеств из нек-рого класса $E$, то условия а) и б) сводятся к условию



\[

\lim _{\varepsilon \downarrow 0} \sup _{A \in E} P\left(A^{\varepsilon} \backslash A^{-\varepsilon}\right)=0

\]



где



\[

A^{\varepsilon}=\mathrm{U}_{x \in A^{\prime}} S_{x, \varepsilon}, A^{-\varepsilon}=S \backslash(S \backslash A)^{\varepsilon}

\]



(когда всякий открытый шар в $S$ связен, $A^{\varepsilon} \backslash A^{-\varepsilon}=$ $\left.=(\partial A)^{\varepsilon}\right)$. Если $S=\mathbb{R} k$ и распределение $P$ абсолютно непрсрывно по мере Лебега, то $P_{n} \Rightarrow P$ тогда и только тогда, когда $P_{n}(A) \rightarrow P(A)$ равномерно по всем борелевским выпуклым множествам $A$.



Пусть $P_{n}, P$ - распределения на метрич. пространстве $S, P_{n} \Rightarrow P$ и $h-$ непрерывное $P$-почти всгоду измеримое отображение $S$ в метрич. пространство $S^{\prime}$; тогда $P_{n} h^{-1} \Rightarrow P h^{-1}$, где для любого распределенпя $Q$ на $S$ распределение $Q h^{-1}$ есть его $h$-образ на $S^{\prime}$ :



\[

Q h^{-1}(A)=Q\left(h^{-1}(A)\right)

\]



для любого борелевского $A \in S^{\prime}$.



Семейство расіределений $\mathscr{P}$ на $S$ наз. с л а 6 о о т нос и т е л в н о к о м п а т т ы м, если всякая последовательность его элементов содержит слабо сходящуюся подпоследовательность. Условие слабой относительной компактности дает теорема Прохорова. Семейство $\mathscr{T}$ наз. п л о т н ы м, если $\forall \varepsilon>0$ существует компакт $K \subset S$ такой, что $P(k)>1-\varepsilon \quad \forall P \in \mathscr{\mathcal { P }}$. Т е ор е м а П р ох о ро в а: если $\mathscr{P}$ плотно, то оно относительно компактно, а если $S$ сепарабельно и полно, то слабая относительная компактность $\mathscr{\mathcal { J }}$ влечет его плотность. В случае, когда $S=\mathbb{R}^{k}$, семейство распределений $\mathscr{P}$ слабо относительно компактно тогда и только тогда, когда соответствующее $\mathscr{P}$ семейство характеристич. функций равностепенно непрерывно в нуле.



Пусть теперь $P_{n}, P-$ распределения на измеримом пространстве $(X, A)$, где $A$ есть $\sigma$-алгебра. Под с х оди м о с т ь ю ІІ о в а р и а ц и и $P_{n}$ к $P$ понимают равномерную сходимость по всем множествам из $A$ или, что равносильно, стремление вариации



\[

\left|P_{n}-P\right|=\left(P_{n}-P\right)^{+}+\left(P_{n}-P\right)^{-}

\]



к нулю; здесь $\left(P_{n}-P\right)^{+}$и $\left(P_{n}-P\right)^{-}$- компоненты разложения Жордана - Хана обобщенной меры $P_{n}-P$.



Jит.: [1] Б и л л и н г с л и П., Сходимость вероятностных мер, пер. с англ., М., 1977; [2] о о ә в М., Теория вероятностей, г а Р а о, Аппроксимация нормальным распределением и асимптотические разложения, пер. англ., 1982. В. В. Сазонов.



РАСПРЕДЕЛЕНИЙ ТИП - совокупность распределений вероятностей случайных величин, получаемых одна из другой каким-либо линейным преобразованием. Точное определение в одномерном случае таково: распределения вероятностей случайных величин $X_{1}$ и $X_{2}$ называют о д н о т и п н ы и, если существуют постоянныс $A$ и $B>0$ такие, тто распределения величин $X_{2}$ и $B X_{1}+A$ совпадают. Соответствующие функции распределения связаны при этом соотношением



\[

F_{2}(x)=F_{1}\left(\frac{x-A}{B}\right)=F_{1}(b x+a),

\]



где $b=1 / B$ и $a=-A / B$.



Таким образом, множество функций распределения разбивается на попарно непересекающиеся типы. При этом, напр., все нормальные распределения образуют один тиі, все равномерные распределения также образуют один тип.



Понятие типа широко используется в предельных теоремах теории вероятностей. Распределение сумм $S_{n}$ независимых случайных величин часто «неограниченно расплываются» при $n \rightarrow \infty$ и сходимость к предельному распределению (например, к нормальному) оказывается возможной только после линейной «нормировки», т. е. для сумм $\frac{s_{n}-a_{n}}{b_{n}}$, где $a_{n}$ и $b_{n}>0$ - нек-рые константы, $b_{n} \rightarrow \infty$ при $n \rightarrow \infty$. При әтом если для каких-либо случайных величин $X_{n}$ распределения величин $\frac{x_{n}-a_{n}}{b_{n}}$ и $\frac{x_{n}-a_{n}^{\prime}}{b_{n}^{\prime}}$ сходятся к невырожденным предельным распроделениям. то эти последние обязательно однотипны. Поэтому можно дать следующее определе-





\end{document}